% 专题十四：数形结合——函数图象与方程根的分布
\section{专题十四：数形结合——函数图象与方程根的分布}

\begin{learnobjective}
\begin{itemize}
  \item 会把方程 $f(x)=g(x)$ 的解看成两个函数图象的交点横坐标；
  \item 会用二次函数的开口、对称轴、顶点、零点画出简图，并据此解不等式；
  \item 会画 $y=|f(x)|$ 的图象（下方翻折上去），据此讨论 $|f(x)|=k$ 的根的个数；
  \item 会用零点定理判断根的存在与区间，并用单调性说明根的个数；
  \item 会用端点函数值的符号判断二次方程两根的位置（根的分布）。
\end{itemize}
\end{learnobjective}

\begin{preview}
\subsection{导语：把方程画出来}
方程的解、不等式的解集，都可以\key{画}出来：
方程 $f(x)=g(x)$ 的解就是两个图象交点的横坐标；
不等式 $f(x)<0$ 的解集就是 $f$ 的图象在 $x$ 轴下方的部分对应的 $x$ 范围。
本专题用五个视角把三个“种子问题”各做一遍，
重点练习「看图说话」和「用端点值定位」。

\subsection{种子问题}
先自己做，再看后面的五个视角。
\begin{enumerate}
  \item 作出 $y=x^2-2x-3$ 的简图（列表关键点即可），
        并据此写出不等式 $x^2-2x-3<0$ 的解集。
  \item 讨论方程 $|x^2-1|=k$ 的实数根的个数（$k$ 为参数）。
  \item 判断方程 $2^x=2-x$ 的实数根的个数，并指出它所在的区间。
\end{enumerate}
\end{preview}

\begin{explore}

\subsection{视角一：来源——二次函数的图象三要素}

\key{要点}
\begin{itemize}
  \item 画二次函数简图抓四点：开口方向、对称轴、顶点、与 $x$ 轴的交点。
  \item 开口向上时，图象在两根之间位于 $x$ 轴下方，两根之外位于 $x$ 轴上方；
        开口向下则相反。
  \item 于是 $f(x)<0$（$f(x)>0$）的解集直接从图上读出，不必死记求根公式。
  \item 顶点纵坐标 $<0$ 且开口向上 $\iff$ 必有两交点；$=0$ 时相切。
\end{itemize}

\begin{example}
作出 $y=x^2-4x+3$ 的简图要点，并写出不等式 $x^2-4x+3<0$ 的解集。
\end{example}
\begin{solution}
$y=x^2-4x+3=(x-1)(x-3)$。开口向上，对称轴 $x=2$，顶点 $(2,-1)$，
与 $x$ 轴交于 $(1,0)$、$(3,0)$。
开口向上时，两根之间图象在 $x$ 轴下方，故 $x^2-4x+3<0$ 的解集为 $(1,3)$。
\key{检验}：取 $x=0$，$0-0+3=3>0$（在两根之外，应为正）\checkmark；
取 $x=2$，$4-8+3=-1<0$（在两根之间，应为负）\checkmark；
取 $x=4$，$16-16+3=3>0$ \checkmark。三个位置的符号都与图象一致。
\key{（回头看：开口方向决定“中间”与“两边”谁正谁负）}
\end{solution}

\begin{problem}
求函数 $y=-x^2+2x+3$ 与 $x$ 轴的交点，并写出不等式 $-x^2+2x+3>0$ 的解集。
\end{problem}
\begin{solution}
$-x^2+2x+3=-(x^2-2x-3)=-(x-3)(x+1)$。
开口向下，对称轴 $x=1$，顶点 $(1,4)$，与 $x$ 轴交于 $(-1,0)$、$(3,0)$。
开口向下时，两根之间图象在 $x$ 轴上方，故 $-x^2+2x+3>0$ 的解集为 $(-1,3)$。
\key{检验}：$x=0$ 时 $y=3>0$ \checkmark（在解集内）；$x=-2$ 时 $y=-4-4+3=-5<0$ \checkmark（在解集外）；
$x=4$ 时 $y=-16+8+3=-5<0$ \checkmark。
\key{（回头看：先定开口，再读上下方，符号一颠倒就全错）}
\end{solution}

\subsection{视角二：变换——绝对值翻折与交点个数}

\key{要点}
\begin{itemize}
  \item $y=|f(x)|$ 的画法：$f(x)\geq0$ 的部分不动；$f(x)<0$ 的部分沿 $x$ 轴\key{翻折}到上方。
  \item 方程 $|f(x)|=k$ 的实根个数 = 图象 $y=|f(x)|$ 与水平直线 $y=k$ 的交点个数。
  \item 关键位置是「谷底」：原函数的最小值 $-h$ 翻折后变成尖峰 $h$，
        $k=h$ 时正好切到尖点，交点个数发生变化。
  \item 讨论时按 $k<0$、$k=0$、$0<k<h$、$k=h$、$k>h$ 分类。
\end{itemize}

\begin{example}
讨论方程 $|x^2-1|=k$ 的实数根的个数。
\end{example}
\begin{solution}
令 $f(x)=x^2-1$，零点 $x=\pm1$，顶点 $(0,-1)$。
$y=|x^2-1|$：在 $(-\infty,-1]$、$[1,+\infty)$ 上与 $x^2-1$ 相同；
在 $(-1,1)$ 上是 $-(x^2-1)=1-x^2$，原谷底 $(0,-1)$ 翻成尖峰 $(0,1)$。
直线 $y=k$ 与其交点个数分类如下：
\begin{itemize}
  \item $k<0$：$0$ 个（$|x^2-1|\geq0$）；
  \item $k=0$：$2$ 个（$x=\pm1$）；
  \item $0<k<1$：$4$ 个（左右两侧各两个，$x=\pm\sqrt{1+k}$、$x=\pm\sqrt{1-k}$）；
  \item $k=1$：$3$ 个（$x=0$ 及 $x=\pm\sqrt2$）；
  \item $k>1$：$2$ 个（只有 $x=\pm\sqrt{1+k}$）。
\end{itemize}
\key{检验}：$k=1$ 时解 $x^2-1=\pm1$，得 $x=0,\pm\sqrt2$，共 $3$ 个 \checkmark；
$k=2$ 时解 $x^2-1=\pm2$，得 $x=\pm\sqrt3$（$x^2=-1$ 舍去），共 $2$ 个 \checkmark。
\key{（回头看：种子问题 2。翻折之后，$k$ 与尖峰高度比较是唯一的分界标准）}
\end{solution}

\begin{problem}
讨论方程 $|x^2-2x|=k$ 的实数根的个数。
\end{problem}
\begin{solution}
令 $f(x)=x^2-2x=x(x-2)$，零点 $x=0,2$，顶点 $(1,-1)$。
$y=|x^2-2x|$：在 $(-\infty,0]$、$[2,+\infty)$ 上不动；
在 $(0,2)$ 上为 $-x^2+2x$，原谷底 $(1,-1)$ 翻成尖峰 $(1,1)$。
分类：
\begin{itemize}
  \item $k<0$：$0$ 个；
  \item $k=0$：$2$ 个（$x=0,2$）；
  \item $0<k<1$：$4$ 个（两侧各两个）；
  \item $k=1$：$3$ 个（$x=1$ 及 $x=1\pm\sqrt2$）；
  \item $k>1$：$2$ 个（只有 $x=1\pm\sqrt{1+k}$）。
\end{itemize}
\key{检验}：$k=1$ 时解 $x^2-2x=\pm1$，$x^2-2x-1=0$ 给 $x=1\pm\sqrt2$，
$x^2-2x+1=0$ 给 $x=1$（重根），共 $3$ 个不同实根 \checkmark；
$k=2$ 时 $x^2-2x-2=0$ 给 $x=1\pm\sqrt3$，$x^2-2x+2=0$ 判别式 $4-8<0$ 无解，
共 $2$ 个 \checkmark。
\key{（回头看：尖峰的高度 $|{-}1|=1$ 来自原函数最小值，先找它再分类）}
\end{solution}

\subsection{视角三：指对——曲线与直线的交点}

\key{要点}
\begin{itemize}
  \item 方程 $f(x)=g(x)$ 的解 = 图象 $y=f(x)$ 与 $y=g(x)$ 的交点横坐标。
  \item 零点定理：若 $F$ 在 $[a,b]$ 上连续且 $F(a)F(b)<0$，则 $F$ 在 $(a,b)$ 内有零点。
  \item 单调的函数至多有一个零点：两个增函数之和仍是增函数，
        增函数与减函数之差仍是增函数。
  \item 指数函数 $y=a^x$（$a>1$）与对数函数 $y=\log_a x$ 都是增函数，常与直线配合使用。
\end{itemize}

\begin{example}
判断方程 $2^x=2-x$ 的实数根的个数与所在区间。
\end{example}
\begin{solution}
法一（数形结合）：$y=2^x$ 是下凸的增函数，$y=2-x$ 是递减直线。
在 $x=0$ 处 $2^0=1<2$，在 $x=1$ 处 $2^1=2>1$，两曲线必相交一次；
又一直增、一直减，交点不可能再有第二个。
法二（代数验证）：令 $F(x)=2^x+x-2$，它在 $\mathbb{R}$ 上连续。
$F(0)=-1<0$，$F(1)=1>0$，由零点定理，存在 $c\in(0,1)$ 使 $F(c)=0$。
又 $2^x$ 与 $x$ 都是增函数，故 $F$ 递增，零点唯一。
\key{检验}：$F\left(\dfrac12\right)=\sqrt2-\dfrac32\approx-0.09<0$，
$F\left(\dfrac34\right)=2^{3/4}-\dfrac54\approx1.68-1.25>0$，
根确实落在 $(0,1)$ 内（约为 $0.54$）\checkmark。
\key{（回头看：种子问题 3）}
\end{solution}

\begin{problem}
判断方程 $\log_2 x=2-x$ 的实数根的个数与所在区间。
\end{problem}
\begin{solution}
令 $G(x)=\log_2x+x-2$（$x>0$）。$G$ 在 $(0,+\infty)$ 上连续，
$G(1)=-1<0$，$G(2)=1>0$，由零点定理存在 $c\in(1,2)$ 使 $G(c)=0$。
又 $\log_2x$ 与 $x$ 在定义域内都是增函数，故 $G$ 递增，零点唯一。
\key{检验}：$G\left(\dfrac32\right)=\log_2\dfrac32-\dfrac12\approx0.585-0.5=0.085>0$，
故根落在 $\left(1,\dfrac32\right)$ 内（约为 $1.46$）\checkmark。
\key{（回头看：定义域 $x>0$ 是先决条件，单调性保证只有一个根）}
\end{solution}

\subsection{视角四：分布——用端点函数值定位根}

\key{要点}
\begin{itemize}
  \item 开口向上的抛物线 $f(x)=ax^2+bx+c$（$a>0$）：
        $f(t)<0\iff t$ 位于两根之间；$f(t)>0\iff t$ 在两根之外。
  \item 两根分居某点 $t$ 两侧 $\iff f(t)<0$（同时要 $\Delta>0$ 保证有两根）。
  \item 两根之积 $x_1x_2=\dfrac ca$、两根之和 $x_1+x_2=-\dfrac ba$ 可用来判断符号与大小关系。
  \item 这类「根的分布」问题，先用端点值列条件，再代回因式分解后的式子检验。
\end{itemize}

\begin{example}
若关于 $x$ 的方程 $x^2-2x+a=0$ 的一个根大于 $2$，另一个根小于 $0$，求 $a$ 的取值范围。
\end{example}
\begin{solution}
$f(x)=x^2-2x+a=(x-1)^2+a-1$，开口向上，判别式 $\Delta=4-4a=4(1-a)$。
方程可写为 $(x-1)^2=1-a$，两根为 $x=1\pm\sqrt{1-a}$（需 $1-a\geq0$）。
小根 $1-\sqrt{1-a}<0\iff\sqrt{1-a}>1\iff a<0$；
大根 $1+\sqrt{1-a}>2\iff\sqrt{1-a}>1\iff a<0$。
两者同时成立当且仅当 $a<0$，此时 $\Delta=4(1-a)>4>0$ 自动满足。
（端点值看法：$f(0)=a<0$ 且 $f(2)=a<0$，说明 $0$ 和 $2$ 都在两根之间，
即小根 $<0<2<$ 大根。）
\key{检验}：取 $a=-1$，方程为 $x^2-2x-1=0$，根 $1\pm\sqrt2\approx-0.41,\ 2.41$，
一个大于 $2$、一个小于 $0$ \checkmark；取 $a=1$，$\Delta=0$，只有一根 $x=1$，不合 \checkmark。
故 $a$ 的取值范围为 $(-\infty,0)$。
\end{solution}

\begin{problem}
若关于 $x$ 的方程 $x^2+bx+1=0$ 的一个根在 $(0,1)$ 内，另一个根大于 $1$，
求 $b$ 的取值范围。
\end{problem}
\begin{solution}
$f(x)=x^2+bx+1$ 开口向上，判别式 $\Delta=b^2-4$，两根之积 $x_1x_2=1>0$，
故两根同号。若一根在 $(0,1)$ 内，记为 $x_1$，则另一根 $x_2=\dfrac1{x_1}>1$，
所以「另一根大于 $1$」自动成立，只需保证一根落在 $(0,1)$ 内。
因为 $f(0)=1>0$，$0$ 在两根之外；要让 $1$ 位于两根之间，只需
\[
  f(1)=1+b+1=b+2<0\quad\Longleftrightarrow\quad b<-2.
\]
此时 $\Delta=b^2-4>4-4=0$，确有两个不等实根。
\key{检验}：取 $b=-3$，方程为 $x^2-3x+1=0$，根 $\dfrac{3\pm\sqrt5}{2}\approx0.38,\ 2.62$，
一个在 $(0,1)$ 内、一个大于 $1$ \checkmark；取 $b=-1$，$\Delta=-3<0$ 无实根，不合 \checkmark。
故 $b$ 的取值范围为 $(-\infty,-2)$。
\key{（回头看：韦达定理先看乘积的符号，能省掉一半讨论）}
\end{solution}

\subsection{视角五：综合——分段讨论与不同类曲线的交点}

\key{要点}
\begin{itemize}
  \item 含 $|x|$ 的方程：令 $t=|x|\geq0$ 化为关于 $t$ 的方程，解出 $t$ 后别忘了 $t$ 非负、
        且 $|x|=t$ 通常给出两个 $x$。
  \item 直线与反比例函数 $y=\dfrac{k}{x}$ 的交点：联立消元成二次方程，
        用判别式判个数，用韦达定理判交点落在哪一支。
  \item 多个知识点混用时，先分类（按绝对值、按参数），再逐段处理。
\end{itemize}

\begin{example}
解方程 $x^2-2|x|-3=0$。
\end{example}
\begin{solution}
左边只含 $|x|$：令 $t=|x|$，则 $t\geq0$，原方程化为
\[
  t^2-2t-3=0\quad\Longleftrightarrow\quad(t-3)(t+1)=0,
\]
得 $t=3$ 或 $t=-1$（舍去，因为 $t\geq0$）。
由 $|x|=3$ 得 $x=3$ 或 $x=-3$。
\key{检验}：$x=3$ 时 $9-6-3=0$ \checkmark；$x=-3$ 时 $9-6-3=0$ \checkmark。
故方程的实数解为 $x=\pm3$。
\key{（回头看：$|x|$ 的方程先换元，解完记得回代并检验非负性）}
\end{solution}

\begin{problem}
讨论直线 $y=x+m$ 与反比例函数 $y=\dfrac3x$ 的图象的交点个数（$m$ 为参数）。
\end{problem}
\begin{solution}
联立 $x+m=\dfrac3x$（$x\neq0$），两边乘以 $x$ 得
\[
  x^2+mx-3=0.
\]
判别式 $\Delta=m^2+12$，对任意实数 $m$ 都有 $\Delta>0$，故总有两个不等实根；
两根之积 $x_1x_2=-3<0$，故两根异号。
而 $y=\dfrac3x$ 的图象在 $x>0$、$x<0$ 各有一支，
所以两个交点分别落在右支和左支上。
\key{检验}：取 $m=0$，交点满足 $x^2=3$，得 $(\sqrt3,\sqrt3)$、$(-\sqrt3,-\sqrt3)$，
确实分属两支、共两个 \checkmark；取 $m=1$，$x^2+x-3=0$ 的两根 $\dfrac{-1\pm\sqrt{13}}2$
一正一负，仍为两个交点 \checkmark。
故对任意实数 $m$，直线与反比例函数图象恰有两个交点。
\key{（回头看：判别式定「有没有」，韦达定理定「在哪里」）}
\end{solution}
\end{explore}

\begin{inclasstest}
\subsection{变式训练}
只改一个条件，重新做一遍，记得代值检验。
\begin{enumerate}
  \item 作出 $y=x^2-2x-8$ 的简图要点，并写出不等式 $x^2-2x-8<0$ 的解集。
  \item 讨论方程 $|x^2-2x-3|=m$ 的实数根的个数。
  \item 判断方程 $3^x=3-x$ 的实数根的个数与所在区间。
  \item 若方程 $x^2+(m-3)x+2=0$ 的两个实数根都在 $(0,2)$ 内，求 $m$ 的取值范围。
  \item 讨论直线 $y=x+m$ 与函数 $y=\dfrac6x$ 的图象的交点个数。
  \item 判断方程 $\log_3x=2-x$ 的实数根的个数与所在区间。
\end{enumerate}
\end{inclasstest}

\begin{afterclass}
\subsection{知识来源对照表}
\begin{center}
\footnotesize
\setlength{\tabcolsep}{4pt}
\renewcommand{\arraystretch}{1.35}
\begin{tabular}{p{3.4cm}p{4.0cm}p{5.6cm}}
\toprule
结论 / 方法 & 最早从哪来 & 后面从哪个角度看 \\
\midrule
二次函数图象三要素 & 视角一（来源） & 视角四（分布） \\
绝对值翻折 & 视角二（变换） & 视角五（综合） \\
零点定理与单调性 & 视角三（指对） & 视角四（分布） \\
端点值定位根的分布 & 视角四（分布） & 视角一（来源） \\
韦达定理判符号 & 视角四（分布） & 视角五（综合） \\
换元解 $|x|$ 方程 & 视角五（综合） & 视角二（变换） \\
\bottomrule
\end{tabular}
\end{center}

\key{下一步}：数形结合的关键是「双向翻译」——
把方程写成两个函数相等，把不等式写成图象在 $x$ 轴的哪一侧，把根的位置写成端点函数值的符号。
同一道题既能画图也能计算，两条路互相印证，就不容易漏解或错判符号。
\end{afterclass}
